% ==============================================================================
%  4.5  Outliers y puntos de influencia
% ==============================================================================
\section{Outliers y puntos de influencia}
\label{sec:t4-outliers}

En algunos problemas, la respuesta observada en unos pocos casos no parece seguir el modelo que
sí se ajusta bien a la gran mayoría de los datos.

\begin{definicion}[Outlier en regresión][def:t4-outlier]
  Un \textbf{outlier en regresión} es un punto que no sigue el patrón lineal general del resto
  de los datos, es decir, que no se ajusta bien al modelo propuesto: su valor ajustado,
  $\est{Y}_i$, está lejos del verdadero valor observado, $Y_i$. Por tanto, estos puntos se
  caracterizan por tener residuos $e_i=Y_i-\est{Y}_i$ altos.
\end{definicion}

\begin{itemize}
  \item El método de mínimos cuadrados es muy sensible a la presencia de outliers, que pueden
    modificar sustancialmente la estimación.
  \item Puede haber outliers tanto en la variable $X$ como en la $Y$.
  \item También puede ocurrir que una observación no sea un outlier en ninguna de las dos
    variables, pero sí en la relación entre $X$ e $Y$ (\cref{sec:t1-descriptivo}).
\end{itemize}
En los modelos de regresión simple es relativamente fácil identificar este tipo de
observaciones mediante diagramas de dispersión.

\begin{figure}[H]
  \centering
  \pgfplotsset{outl/.style={mini, width=3.5cm, height=3.4cm, xmin=0, xmax=15, ymin=0, ymax=15}}
  \def\nube{(1,1) (1,2) (2,2) (2,3) (2,4) (3,3) (3,4) (4,3) (4,4) (4,5) (5,5) (5,6) (6,5)
    (6,6) (6,7) (7,7) (7,8) (8,8) (8,9) (8,10)}
  \begin{subfigure}[t]{0.24\linewidth}
    \centering
    \begin{tikzpicture}
      \begin{axis}[outl]
        \addplot[puntos, mark size=1.2pt] coordinates {\nube};
        \addplot[only marks, mark=*, mark size=1.8pt, colRes] coordinates {(13,5)};
      \end{axis}
    \end{tikzpicture}
    \caption{Valor extremo de $X$.}
  \end{subfigure}\hfill
  \begin{subfigure}[t]{0.24\linewidth}
    \centering
    \begin{tikzpicture}
      \begin{axis}[outl]
        \addplot[puntos, mark size=1.2pt] coordinates {\nube};
        \addplot[only marks, mark=*, mark size=1.8pt, colRes] coordinates {(5,14)};
      \end{axis}
    \end{tikzpicture}
    \caption{Valor extremo de $Y$.}
  \end{subfigure}\hfill
  \begin{subfigure}[t]{0.24\linewidth}
    \centering
    \begin{tikzpicture}
      \begin{axis}[outl]
        \addplot[puntos, mark size=1.2pt] coordinates {\nube};
        \addplot[only marks, mark=*, mark size=1.8pt, colRes] coordinates {(13,14)};
      \end{axis}
    \end{tikzpicture}
    \caption{Extremo en $X$ y en $Y$.}
  \end{subfigure}\hfill
  \begin{subfigure}[t]{0.24\linewidth}
    \centering
    \begin{tikzpicture}
      \begin{axis}[outl]
        \addplot[puntos, mark size=1.2pt] coordinates {\nube};
        \addplot[only marks, mark=*, mark size=1.8pt, colRes] coordinates {(1,10)};
      \end{axis}
    \end{tikzpicture}
    \caption{Alejado de la relación.}
  \end{subfigure}
  \caption{Tipos de outliers (en rojo) en un diagrama de dispersión.}
\end{figure}

Hay, por tanto, dos tipos de atipicidad:
\begin{itemize}
  \item \textbf{outliers univariantes}, atípicos en $X$ o en $Y$;
  \item \textbf{outliers en la relación}, que no tienen por qué ser atípicos en ninguna de las
    dos variables.
\end{itemize}

\begin{definicion}[Punto de influencia][def:t4-influencia]
  Un \textbf{punto de influencia} es un outlier que tiene un efecto muy importante en la recta de
  regresión ajustada.
\end{definicion}

\begin{figure}[H]
  \centering
  \begin{subfigure}[t]{0.45\linewidth}
    \centering
    \begin{tikzpicture}
      \begin{axis}[mini, xmin=-0.1, xmax=2.1, ymin=-0.8, ymax=3.6]
        \addplot[puntos, mark size=1.2pt] table {datos/tema4/infl_sin.dat};
        \addplot[only marks, mark=triangle*, mark size=2.5pt, colRes] coordinates {(2,0)};
        \addplot[recta] table[y={create col/linear regression={y=y}}] {datos/tema4/infl_con.dat};
      \end{axis}
    \end{tikzpicture}
    \caption{Con el punto de influencia (triángulo rojo).}
  \end{subfigure}\hfill
  \begin{subfigure}[t]{0.45\linewidth}
    \centering
    \begin{tikzpicture}
      \begin{axis}[mini, xmin=-0.1, xmax=2.1, ymin=-0.8, ymax=3.6]
        \addplot[puntos, mark size=1.2pt] table {datos/tema4/infl_sin.dat};
        \addplot[recta] table[y={create col/linear regression={y=y}}] {datos/tema4/infl_sin.dat};
      \end{axis}
    \end{tikzpicture}
    \caption{Sin el punto de influencia.}
  \end{subfigure}
  \caption{Recta ajustada con y sin un punto de influencia (datos simulados).}
\end{figure}

Si se encuentran puntos de influencia en un conjunto de datos, hay que tener en cuenta que pueden
ser observaciones que no representan el comportamiento de interés, y comparar las decisiones que
se toman cuando se incluyen y cuando no.

\subsection{Outliers en \texorpdfstring{$X$}{X}: el leverage}
\label{sec:t4-outliers-x}

Los outliers en $X$ pueden identificarse por valores grandes del leverage $h_{ii}$
(\cref{def:t4-leverage}).
\begin{itemize}
  \item En general, $0\le h_{ii}\le1$ y $\sumi h_{ii}=2$ (\cref{prop:t4-prop-leverage}).
  \item Valores grandes del leverage indican que el valor $x_i$ está alejado del centro de todas
    las observaciones: $(x_i-\media{X})^2$ grande.
  \item Se considera que un leverage es grande si es mayor que dos veces la media de los
    leverages:
    \[
      h_{ii} > 2\media{h} = \frac2n\sumi h_{ii} = \frac4n .
    \]
\end{itemize}

Los puntos con leverage grande son \textbf{potenciales} puntos de influencia, es decir, pueden
condicionar de manera importante el ajuste del modelo:
\begin{itemize}
  \item como $\Var(e_i)=\sigma^2(1-h_{ii})$, cuanto mayor es $h_{ii}$ menor es la varianza de
    $e_i$;
  \item puesto que $0\le h_{ii}\le1$, cuanto más próximo a 1 es $h_{ii}$, más próxima a cero es la
    varianza del residuo de la observación $i$-ésima;
  \item por tanto, en las observaciones con $h_{ii}$ grande, $\est{Y}_i$ tiende a estar cerca del
    valor observado $Y_i$;
  \item en el caso extremo e hipotético $h_{ii}=1$, la recta ajustada estaría forzada a pasar por
    el valor observado $(x_i,Y_i)$.
\end{itemize}

\subsection{Outliers en \texorpdfstring{$Y$}{Y}: los residuos}
\label{sec:t4-outliers-y}

Los residuos permiten identificar los outliers en la variable respuesta: valores altos o bajos
indican observaciones que se ajustan mal al modelo. Los residuos regulares, $e_i=Y_i-\est{Y}_i$,
no tienen una escala conocida, por lo que es difícil definir qué es un ``valor extremo''; por
eso siempre resulta más útil utilizar los residuos estudentizados, $r_i$
(\cref{def:t4-estudentizados}), o los estudentizados eliminados, $t_{(i)}$
(\cref{def:t4-eliminado-estudentizado}).
\begin{itemize}
  \item Tanto $r_i$ como $t_{(i)}$ se distribuyen según una $t$ de Student (aproximadamente en el
    caso de $r_i$). Para un tamaño $n$ razonable, un valor de $|r_i|$ o de $|t_{(i)}|$
    \textbf{mayor que 3} indica que la observación es un outlier.
  \item Los valores de $|r_i|$ o $|t_{(i)}|$ \textbf{entre 2 y 3} deben considerarse
    sospechosos.
  \item Muchas observaciones con residuos extremos podrían indicar que el modelo no es adecuado.
  \item Puede ser más útil utilizar $t_{(i)}$ que $r_i$, ya que la observación no interviene en
    el cálculo de su propio residuo.
\end{itemize}

\paragraph{Contraste de outliers.} De manera más formal, para cada observación puede plantearse
el contraste
\[
  H_0\colon \text{la observación } i \text{ no es un outlier}
  \qquad\text{frente a}\qquad
  H_1\colon \text{la observación } i \text{ es un outlier},
\]
donde $H_0$ significa que la observación $i$ sigue el modelo ajustado. Bajo $H_0$ (y con errores
normales), $t_{(i)}\sim\tdist{n-3}$, de modo que se rechaza $H_0$ a nivel $\alpha$ si
$|t_{(i)}|>\tq{n-3}{1-\alpha/2}$. Como se contrastan los $n$ residuos, aparece el problema de las
comparaciones múltiples (\cref{sec:t3-bonferroni}), y se aplica el ajuste de Bonferroni con
$\alpha^*=\alpha/n$.

\begin{resumen}[Identificación de outliers en $Y$]
  Se contrasta cada uno de los $n$ residuos para determinar si corresponde a un outlier: la
  observación $i$ se considera un outlier si
  \[
    |t_{(i)}| > \tq{n-3}{1-\alpha/(2n)} .
  \]
\end{resumen}

\subsection{Medidas de influencia}

Un outlier no siempre es un punto de influencia, ni los outliers son los únicos puntos
sospechosos de influir demasiado en el ajuste.
\begin{itemize}
  \item Un outlier en la variable $Y$ puede ser influyente. Serán potencialmente más influyentes
    los puntos con residuos mucho mayores (o menores) que el resto de residuos de la muestra.
  \item Los puntos con un valor particularmente alto o bajo de la variable $X$ pueden tener una
    gran influencia en la estimación de los dos coeficientes. Estos puntos son potencialmente
    influyentes.
\end{itemize}
Así, con los residuos estudentizados (eliminados) y los leverages se detectan los puntos
potencialmente influyentes, pero es necesario evaluar su \textbf{influencia efectiva}. Para ello
se utilizan las siguientes medidas:
\begin{itemize}
  \item la \textbf{distancia de Cook}, que mide la influencia de una observación en todos los
    valores ajustados;
  \item los \textbf{DFFits}, que miden la influencia de una observación en su propio valor
    ajustado;
  \item los \textbf{DFBeta}, que miden la influencia de una observación en un coeficiente de
    regresión particular.
\end{itemize}
Las tres se basan en comparar el ajuste con todas las observaciones y el ajuste sin la
observación $i$-ésima (\cref{sec:t4-eliminados}).

\paragraph{Distancia de Cook.} Determina la influencia de una observación en todos los valores
predichos; valores grandes indican que la observación tiene mucha influencia. Mide cómo cambian
los estimadores de los coeficientes de regresión cuando se elimina una observación, utilizando la
distancia de Mahalanobis para cuantificar el cambio:
\[
  D_i = \frac{\sum_{j=1}^n\big(\est{Y}_j-\est{Y}_{j(i)}\big)^2}{2\cdot MSE}
  = \frac{e_i^2}{2\cdot MSE}\left(\frac{h_{ii}}{(1-h_{ii})^2}\right).
\]
Su valor se compara con los de la distribución $\Fdist{2}{n-2}$: si $D_i$ se sitúa en torno a la
mediana de la $\Fdist{2}{n-2}$ o por encima de ella, la observación tiene una influencia
importante en el ajuste; si queda por debajo de su percentil 10 o 20, apenas influye.

\paragraph{DFFits.} Determina la influencia de una observación en su propia predicción. Es una
medida estandarizada del cambio que sufre $\est{Y}_i$ al eliminar la observación $i$:
\[
  \DFFits_i = \frac{\est{Y}_i-\est{Y}_{i(i)}}{\sqrt{MSE_{(i)}\cdot h_{ii}}} .
\]
En general, una observación se considera influyente si
\(
  |\DFFits_i| > 2\sqrt{2/n} .
\)

\paragraph{DFBeta.} Hay un DFBeta por cada parámetro y cada observación, y determina la influencia
de cada observación en la estimación de cada parámetro. Es una medida estandarizada del cambio que
sufre $\est{\beta}_j$, $j=0,1$, al eliminar la observación $i$:
\[
  \DFBeta_{ji} = \frac{\est{\beta}_j-\est{\beta}_{j(i)}}{\sqrt{\widehat{\Var}\big(\est{\beta}_{j(i)}\big)}},
\]
donde el error estándar del denominador se calcula con $MSE_{(i)}$ en lugar de $MSE$
(\cref{teo:t3-dist-b0,teo:t3-dist-b1}); por ejemplo, para la pendiente es
$\sqrt{MSE_{(i)}/\SSX}$, con $\SSX$ calculado con todas las observaciones. Se consideran altos los valores absolutos mayores que 1 (en muestras
pequeñas o medianas) o que $2/\sqrt{n}$ (en muestras grandes).

\subsection{Tipos de outliers y puntos de influencia}

\begin{center}
  \small
  \renewcommand{\arraystretch}{1.3}
  \begin{tabular}{l>{\raggedright\arraybackslash}p{4.6cm}>{\raggedright\arraybackslash}p{4.6cm}}
    \toprule
    & \textbf{No outlier en $Y$} & \textbf{Outlier en $Y$} \\
    \midrule
    \textbf{No outlier en $X$} &
      Puntos normales (A): residuos pequeños y leverages pequeños. &
      Outlier, problema leve (B): residuos grandes y leverages pequeños. \\
    \addlinespace
    \textbf{Outlier en $X$} &
      Punto de influencia ``bueno'' (C): leverage grande y medidas de influencia efectiva
      pequeñas. &
      Punto de influencia ``malo'' (D): leverage grande y medidas de influencia efectiva
      grandes. \\
    \bottomrule
  \end{tabular}
\end{center}

\begin{figure}[H]
  \centering
  \begin{tikzpicture}
    \begin{axis}[width=7.5cm, height=5.5cm, ticks=none, axis lines=box, xmin=-0.5, xmax=11,
        ymin=-0.5, ymax=11]
      \addplot[black, no marks, domain=0:10.6] {0.15+0.93*x};
      \addplot[puntos, mark size=1.4pt] coordinates {(1,1.2) (1.5,0.7) (1.5,1.7) (2,2.1)
        (2.3,1.4) (2.7,2.6) (3.2,2.1) (3.6,2.8) (3.7,3.6) (4.1,4.3) (4.4,3.2) (4.6,4.1)
        (5,4.9) (5.3,4.3)};
      \addplot[only marks, mark=*, mark size=1.8pt, colRes] coordinates {(2.2,6.6) (10,9.6) (10,3.8)};
      \node[font=\footnotesize, anchor=north] at (axis cs:2.9,0.9) {A};
      \node[font=\footnotesize, anchor=south east] at (axis cs:2.1,6.7) {B};
      \node[font=\footnotesize, anchor=south east] at (axis cs:9.9,9.7) {C};
      \node[font=\footnotesize, anchor=south east] at (axis cs:9.9,3.9) {D};
    \end{axis}
  \end{tikzpicture}
  \caption{A: puntos normales; B: outlier; C: punto de influencia bueno; D: punto de influencia
    malo.}
\end{figure}

\paragraph{Tratamiento.} Se eliminan los casos que condicionan de manera importante el análisis y
se estudian las causas de la aparición de dichas observaciones.
